% Requires: \usepackage{booktabs} \usepackage{threeparttable}
\begin{table}[htbp]
\centering
\footnotesize
\caption{Difference tests against the equal-weight benchmark.}
\label{tab:inference}
\begin{threeparttable}
\begin{tabular}{lrrrr}
\toprule
 & {EPO} & {PPP} & {Proposed (hist.)} & {Proposed (fwd.)} \\
\midrule
\multicolumn{5}{l}{\textit{Panel A: Mean return}} \\
Difference & -3.96 & -0.72 & 18.68$^{***}$ & 22.86$^{***}$ \\
$t$-statistic & -1.48 & -0.40 & 2.73 & 3.19 \\
Bootstrap $p$-value & [0.170] & [0.688] & [0.004] & [0.000] \\
\addlinespace
\multicolumn{5}{l}{\textit{Panel B: Sharpe ratio}} \\
Difference & -0.22 & -0.10 & 0.30 & 0.45 \\
$t$-statistic & -1.27 & -1.04 & 1.16 & 1.64 \\
Bootstrap $p$-value & [0.227] & [0.306] & [0.269] & [0.111] \\
\addlinespace
\multicolumn{5}{l}{\textit{Panel C: CRRA certainty equivalent}} \\
Difference & -3.11 & -3.29 & 2.74 & 8.71 \\
$t$-statistic & -1.03 & -1.27 & 0.36 & 1.05 \\
Bootstrap $p$-value & [0.318] & [0.261] & [0.723] & [0.304] \\
\bottomrule
\end{tabular}
\begin{tablenotes}[flushleft]\footnotesize
\item Sample 2015-02 to 2026-01 (132 monthly observations). Returns are net of 10~bps of one-way transaction cost. Each column reports the strategy minus the equal-weight benchmark, whose levels are in Table~\ref{tab:performance}.
\item All tests are two-sided and paired. Each is a delta method over sample moments with a prewhitened quadratic-spectral HAC long-run covariance (Andrews, 1991; Andrews and Monahan, 1992); the reported $p$-value is from a studentized circular block bootstrap with block length 5 and 4999 resamples (Ledoit and Wolf, 2008).
\item The panels are ordered from the statistic that does not adjust for risk to the one that penalises it most heavily. A mean return difference is not scale-invariant: it rises with leverage without the accompanying volatility being charged for. Significance in Panel A alongside insignificance in Panels B and C is therefore a statement about volatility, not about skill.
\item With three tests across four strategies these are twelve comparisons; the notes on multiple testing in the text apply.
\item $^{***}$, $^{**}$ and $^{*}$ denote significance at the 1\%, 5\% and 10\% levels.
\end{tablenotes}
\end{threeparttable}
\end{table}
